
\documentclass[11pt, oneside]{article}   	% use "amsart" instead of "article" for AMSLaTeX format
\usepackage{geometry}                		% See geometry.pdf to learn the layout options. There are lots.
\geometry{letterpaper}            
\usepackage[parfill]{parskip}    		% Activate to begin paragraphs with an empty line rather than an indent
\usepackage{graphicx}				% Use pdf, png, jpg, or eps§ with pdflatex; use eps in DVI mode
								% TeX will automatically convert eps --> pdf in pdflatex		
\usepackage{amssymb}

%SetFonts

%SetFonts


\title{Special Comment Lines; Hidden Settings}
% \author{The Author}
\date{}							% Activate to display a given date or no date

\begin{document}
\maketitle
%\section{}
%\subsection{}
\section{Magic Comment Lines}
A ``Magic Comment Line'' is a command in the source which is interpreted by TeXShop. Such lines must occur in the first twenty lines of a source file.
%  and begin with the comment symbol \%.
%TeXShop understands the following special comment lines:

\begin{itemize}
\item \% !TEX program =
\item \% !TEX encoding =
\item \% !TEX root =
\item \% !TEX projectfile =
\item	\% !TEX numberingCorrection = 
\item \% !TEX useTabs
\item \% !TEX useTabsWithFiles
\item \% !TEX tabbedFile\{ \}(optional short name)
\item \% !TEX pdfSinglePage
\item \% !TEX spellcheck =
\item \% !BIB program =
\item \% !TEX parameter =
\item \% !TEX useOldSyncParser
\item \% !TEX useConTeXtSyncParser
\item \% !TEX useAlternatePath
\item \% !TEX bookDisplay
\item \% !TEX standardDisplay
\item \% !TEX PageDirectionL2R
\item	\% !TEX PageDirectionR2L

\end{itemize}

The older forms ``\% !TEX TS-program ='' and ``\% !TEX TS-parameter ='' are also accepted.

\section{Program Examples}

TeX source files can be typeset by many programs: TeX, LaTeX, pdfLaTeX, luaLaTeX, etc. TeXShop has an ``Engine mechanism" to extend this list to user defined shell scripts, which can run a sequence of programs one after another. The Special Comment ``program'' line determines which program typesets a particular source file. But using this line is never required because TeXShop has other ways to determine the typesetting program.

If the special comment line is present, it overrides all other methods of determining the typesetting program. Thus if a particular source uses Unicode and requires XeLaTeX, adding a Special Comment line will guarantee that it is always typeset by the correct program.

The syntax of Special Comment lines is very picky about required spaces, cases, etc. Thus I recommend 
Ramon Figueros-Centeno's ``Program'' macro, because it lists all active engines and the user can simply click the required engine; Ramon's macro will automatically insert a Special Comment with exactly the right syntax.

Below are some examples:

\begin{verbatim}
     % !TEX program = tex               
     % !TEX program = latex
     % !TEX program = pdftex
     % !TEX program = pdflatex
     % !TEX program = xelatex
     % !TEX program = pdflatexmk
\end{verbatim}

Note that ``tex'' and ``latex'' produce TeX and DVI typesetting, in which TeX outputs a dvi file, which is then processed to produce a final pdf file. ``pdftex'' and ``pdflatex'' use pdftex to directly output pdf files.
Note also that ``xelatex'' and ``pdflatexmk'' are engines defined by scripts in $\sim$Library/TeXShop/Engines.

\section{Encoding Examples}

These Special Comment lines fix the encoding used to read and write a source file, overriding all other methods of setting this encoding. It is important to use the exact name for an encoding used by the Macintosh, so I again recommend Ramon Figueros-Centeno's ``Encoding'' macro, which lists all possibilities and requires only a click on a choice to produce the correct syntax line.
 
Without such a line, TeXShop detects the encoding when it opens a file: it looks for a byte order mark, the encoding recorded by the program that wrote the file, and valid UTF-8, and otherwise chooses the encoding that reads the file most plausibly among the encoding set in TeXShop Settings and the encodings of your languages. An encoding chosen in the Open dialog is used as it is.

 \begin{verbatim}    
     % !TEX encoding = UTF-8 Unicode
     % !TEX encoding = Iso Latin 9
     % !TEX encoding = Mac Chinese Traditional
 \end{verbatim}    
   
 \section{Tabbed Window Examples}
 
 The examples in this section still work, but they are mostly replaced by the tabbing features introduced in TeXShop 4.72. Read the changes document about that version for details.
  
 The special comment ``root'' is explained elsewhere in the TeXShop Help Panel, and used when a project is divided into a commanding root file and included chapter files. 
 
 macOS supports tabbed windows; see the macOS documentation for creating tabs without TeXShop.
  
 TeXShop has three special comment commands to extend this build-in system support for tabs. 
 
 \begin{verbatim}
     % !TEX useTabs
	
     % !TEX useTabsWithFiles
     % !TEX tabbedFile{Galois.aux}(Aux)
     % !TEX tabbedFile{Galois.log}
     % !TEX tabbedFile{~/Graphics/faireyes.eps}
 \end{verbatim}
 The first line above assumes a root file containing ``include'' lines for the chapter sources. This special comment creates a window containing the root and various chapter sources as tabs in a single window.
 In this case, TeXShop searches the root file for the include-lines which name these chapter files.
 
 The last four lines describe a more general method of creating tabbed views. This method requires more
 information from the user, but produces more flexible tabs. The ``useTabsWithFiles'' line introduces the method, which then produces a tabbed window for each ``tabbedFile'' line. This line contains a partial or full path to the file in curly brackets. The tabs will be named using these paths, but this can create long tab names, so it is possible to provide  shorter names for the tabs inside round brackets. These
 shorter names are optional.
 
 \newpage
 \section{More Examples}
 \begin{verbatim}
       % !TEX numberingCorrection = 0 + current - desired 
 \end{verbatim}
This command lists the number of preliminary pages before the main material is given on pages starting with page 1. 
TeXShop uses this information to synchronize the PageNumber box on the toolbar with actual page numbers in the document.

 \begin{verbatim}
      % !TEX pdfSinglePage
 \end{verbatim}
 This command was proposed by a user for  Beamer. The user preferred Multiple Page mode for the  display of articles and books, but wanted Beamer documents to appear as single pages in preview, ready to be projected onto a screen. Place this special comment in the source of any document whose preview window should appear in single page mode.
 
  \begin{verbatim}
      % !TEX spellcheck = German
 \end{verbatim}
 This command causes the Spelling and Grammar panel to use the German dictionary when the document is active, even if by default another dictionary is generally used.
 
  \begin{verbatim}
      % !BIB program = biber
 \end{verbatim}
 This causes the BibTeX command to run an alternate program, in this case biber.    
  
  \begin{verbatim}
      % !TEX parameter = 
 \end{verbatim}
When engine scripts run, they are passed a parameter giving the full path to the source file. This command passes a second parameter to such engine scripts. The engine script can ignore this parameter, but in some cases it is useful to write a script which is able to perform several tasks, depending on the new parameter.  
     
  \begin{verbatim}
      % !TEX useOldSyncParser
      % !TEX useConTeXtSyncParser
 \end{verbatim}
Synctex was written by Jerome Laurens. It causes a TeX engine to output a .synctex file with information needed to jump from a spot in the source to a corresponding spot in the preview, and vice-versa. Laurens also provides parser code for front ends, making it possible for front end authors to provide syncing ability without much extra work. In TeX Live 2017, Laurens fixed bugs and extended both pieces of code.

The ConTeXt system for luaTeX is written by Hans Hagen. Hagen often replaces pieces of TeX Live with his own code in ConTeXt. In 2017, Hagen wrote his own version of the sync code, but unfortunately he based it on Lauren's earlier version 2016 of the code, The 2017 parser cannot handle this code.

After a good deal of work, TeXShop contains both the 2016 and 2017 versions of the parser. By default it uses the 2017 version, but ConTeXt users can switch to the older parser using the special comment above. I don't yet know what will happen in 2018, but shoehorning in both parser libraries was an unpleasant task I am not likely to attempt again.

(Added in 2021): Hans modified the sync code in ConTeXt, Then he wrote routines which front ends can call to obtain  sync information from the ConTeXt synctex file. In short, Laurens' parser code is completely replaced by code in ConTeXt. The special comment line "useConTeXtSyncParser" causes TeXShop to call this ConTeXt code. From 2021 on, this is the preferred sync method for ConTeXt users. 

One consequence of this development is that Hans can  modify sync code if he wishes without breaking sync in TeXShop. This is a welcome development. Thanks to Nicola Vitacolonna for calling this to my attention and urging TeXShop to adopt the new calls.


\begin{verbatim}
       % !TEX useAlternatePath
\end{verbatim}
 
 Most TeXShop users pair it with the TeX Live distribution of TeX. This distribution is updated once a year, when new versions of the various executable files are released. The style files, class files, fonts, etc. in TeX Live are updated daily.
 
 This organization works for most TeX binaries. ConTeXt is an exception because Hans Hagen updates it regularly and often. Luckily, there is an easy way to install and update ConTeXt, supported by a beautiful series of web pages called the "ConTeXt Garden." See https://wiki.contextgarden.net/Main\_Page. Notice the link on this page titled "Install ConTeXt and start typesetting." This link downloads and installs everything needed to typeset using ConTeXt, placing it in a location of the user's choice. One typical place is $\sim$/bin/context. 
 
TeXShop needs to be reconfigured to use this distribution rather than TeX Live. This is easy for someone who always typesets with ConTeXt, but troublesome for users who write some documents in LaTeX and some in ConTeXt. The Typesetting category of TeXShop Settings has a setting called "Alternate Path" where users can insert the full path to their ConTeXt distribution. The special command line "useAlternatePath" can then be added to the top of a ConTeXt document to use that path to ConTeXt rather than the standard path to TeX Live.

\begin{verbatim}
       % !TEX projectfile = chapter2.tex
\end{verbatim}

When there is no SyncTeX information, TeXShop synchronizes by searching for text. To find the file for a spot in the Preview window, it searches the files named in the include and input commands of the root file. ConTeXt includes files differently, so list the files of such a project with these lines in the root file, one per line. A name is a full path or a path relative to the root file. Unlike the other magic lines, these lines may occur anywhere in the root file.

\newpage
\begin{verbatim}
       % !TEX bookDisplay
       % !TEX standardDisplay
       
       % !TEX PageDirectionL2R
       % !TEX PageDirectionR2L
\end{verbatim}

In ``double page'' and ``double multi-page'' modes, the Preview window shows the first page by itself and then shows a series of double pages. By convention, there are an odd number of pages in a book before the actual content begins, so this causes the left and right pages to appear as they would if a user opened the actual book.  These magic lines allow users to change this behavior on a document-by-document basis. The first magic line causes the first page to appear by itself, and the second causes it to appear as a pair with the second page.

 In Japan, text can be either vertical or horizontal, and vertical text pages should appear from right to left.  The last two lines set this on a document by document basis.


 
\section{Hidden Settings}
This reference is also available from Settings, General, Advanced. Spelling methods and their rules are in Spelling; keyboard input editors are in Keyboard. Disabling a dependent option does not reset its saved value.

TeXShop has a large number of settings which make it possible for users to customize the behavior of the program to their liking. The most important ones are in TeXShop Settings, but more obscure settings are hidden and only available through the Terminal in /Applications/Utilities. To change such a setting, quit TeXShop, enter the command in Terminal and push RETURN, and then restart TeXShop. To return to the default, use for example
\begin{verbatim}
     defaults delete TeXShop FlashDelay
\end{verbatim}

Each command below changes a setting from its default. These are all the hidden settings of this version of TeXShop.

\newpage
{\em Editing}
\begin{description}
\item[\texttt{defaults write TeXShop ColorImmediately NO}] \hfill \\
Do not color the syntax of a long file (over 100,000 characters) while it is being opened.
\item[\texttt{defaults write TeXShop ResetSourceTextColorEachTime YES}] \hfill \\
Set the source text color again whenever a source window becomes active; this helps with Japanese input methods and dark backgrounds.
\item[\texttt{defaults write TeXShop IndexColorStart YES}] \hfill \\
Start with the index coloring tool on.
\item[\texttt{defaults write TeXShop SyntaxColorFootnote NO}] \hfill \\
Do not color the argument of \textbackslash footnote, \textbackslash autocite, \textbackslash footcite and the \textbackslash cite commands.
\item[\texttt{defaults write TeXShop brieflyFlashYellowForMatch NO}] \hfill \\
When a closing bracket is typed and bracket matching is on, briefly select the matching bracket instead of flashing it in yellow.
\item[\texttt{defaults write TeXShop ContiguousTextLayout YES}] \hfill \\
Lay out the whole text of a source window at once instead of only the visible part. Long documents then open and scroll more slowly.
\item[\texttt{defaults write TeXShop RightJustifyIfAny YES}] \hfill \\
When right justification is on, right justify a line if it contains any right-to-left character, rather than only when its first characters are right-to-left.
\item[\texttt{defaults write TeXShop ConvertLF NO}] \hfill \\
Keep carriage returns in the text. By default, line endings in files that are opened and in text that is typed or inserted are changed to line feeds.
\item[\texttt{defaults write TeXShop IncrementalSearch NO}] \hfill \\
Do not search incrementally in the Find Bar. This applies only when the Find Bar is chosen in TeXShop Settings.
\item[\texttt{defaults write TeXShop CommandCompletionAlternateMarkShortcut NO}] \hfill \\
Turn off the shortcuts that move to the next and previous completion marks: Option and Control with the command completion key (Escape or Tab). The value must be written as NO, as shown.
\item[\texttt{defaults write TeXShop MathExpressionCompletion 1}] \hfill \\
Show the results of math expressions while typing (macOS 15 and later). 0 turns this off (the default), 1 turns it on and 2 follows the system setting. It applies to windows opened afterwards.
\end{description}

{\em Spell Checking}
\begin{description}
\item[\texttt{defaults write TeXShop OriginalSpelling YES}] \hfill \\
Use the standard Apple spell checking language handling instead of switching the dictionary for each document.
\item[\texttt{defaults write TeXShop ExtraCommandsNotToCheckParameters -array '\textbackslash mycite'}] \hfill \\
With TeX-aware spell checking off and the option not to check parameters on, do not check the spelling of the parameters of these commands either. Each name includes its backslash.
\item[\texttt{defaults write TeXShop ExceptionListExcludesParameters YES}] \hfill \\
With TeX-aware spell checking off and the option not to check parameters on, check no parameters except those of \textbackslash emph, \textbackslash verbatim and the commands in ExtraCommandsToCheckParameters.
\item[\texttt{defaults write TeXShop ExtraCommandsToCheckParameters -array '\textbackslash mycaption'}] \hfill \\
The commands whose parameters are still checked when ExceptionListExcludesParameters is YES.
\end{description}

{\em Tags and Labels}
\begin{description}
\item[\texttt{defaults write TeXShop TagSections NO}] \hfill \\
Do not list \textbackslash chapter, \textbackslash section and the like in the Tags menu; only the lines marked as tags are listed.
\item[\texttt{defaults write TeXShop ConTeXtTags YES}] \hfill \\
List the ConTeXt section commands, such as \textbackslash title and \textbackslash subject, in the Tags menu instead of the LaTeX ones.
\item[\texttt{defaults write TeXShop CreateTagList NO}] \hfill \\
Leave the Tags menu empty.
\item[\texttt{defaults write TeXShop CreateLabelList NO}] \hfill \\
Leave the Labels menu empty.
\item[\texttt{defaults write TeXShop UseNewTagsAndLabels NO}] \hfill \\
Rebuild the Tags and Labels menus shortly after each change to the text, as older versions did, instead of when the menu is opened.
\end{description}

{\em Typesetting and Files}
\begin{description}
\item[\texttt{defaults write TeXShop EncodingAutoDetection NO}] \hfill \\
Do not detect the encoding of a file without an encoding magic line; open it in the encoding of TeXShop Settings, as older versions did.
\item[\texttt{defaults write TeXShop UseOldHeadingCommands YES}] \hfill \\
Also accept the old first line forms of the magic lines: \%\&pdflatex and similar lines and \%\&program= for the typesetting program, \%\&encoding= for the encoding and \%SourceDoc for the root file.
\item[\texttt{defaults write TeXShop RedConsoleAfterError YES}] \hfill \\
Show text typed in the console to answer a TeX error, and all later output of that run, in red.
\item[\texttt{defaults write TeXShop OtherTeXExtensions -array 'mkiv'}] \hfill \\
Treat files with these extensions as TeX sources, so that they can be typeset and are syntax colored. Give the extensions without the dot.
\item[\texttt{defaults write TeXShop AggressiveTrashAUX YES}] \hfill \\
Make Trash Aux Files always work as when Option is held down: it removes the auxiliary files in the folder of the root file and its subfolders, whatever their names.
\item[\texttt{defaults write TeXShop OtherTrashExtensions -array 'nav' 'snm'}] \hfill \\
Also remove files with these extensions in Trash Aux Files. Give the extensions without the dot.
\end{description}

{\em Preview}

The default page direction and whether to show the first page alone as a cover
can be selected in Settings $\rightarrow$ Preview. The PDF preview toolbar
also provides page layout and two-page arrangement menus for the current document.
Magic comments override PDF reading direction, which overrides the application
setting. A temporary toolbar choice takes precedence until restored to the
document settings or the document is closed.

\begin{description}
\item[\texttt{defaults write TeXShop InterpolationValue 4}] \hfill \\
The image interpolation used to draw the Preview window when antialiasing is on: 0 (system default), 1 (none), 2 (low), 4 (medium) or 3 (high, the default).
\item[\texttt{defaults write TeXShop FlashFix NO}] \hfill \\
Do not cover the Preview window briefly while a new version of the PDF is loaded; the cover prevents flashing.
\item[\texttt{defaults write TeXShop FlashDelay 0.25}] \hfill \\
How long, in seconds, the cover stays; the default is 0.5.
\item[\texttt{defaults write TeXShop UseOutline NO}] \hfill \\
Do not build the PDF outline in the Preview drawer.
\item[\texttt{defaults write TeXShop PreviewDrawerOpen YES}] \hfill \\
Open the Preview drawer when a Preview window first appears.
\item[\texttt{defaults write TeXShop LinkPopups NO}] \hfill \\
Start each Preview window with link popups off. The Preview menu turns them on and off for a window.
\item[\texttt{defaults write TeXShop LeftRightArrowsAlwaysPage YES}] \hfill \\
Make the left and right arrow keys always move to the previous or next page, even when the page can be scrolled sideways.
\item[\texttt{defaults write TeXShop GlassMaxMagnification 400}] \hfill \\
The magnifying glass works only while the Preview is shown at no more than this magnification, in percent; the default is 250.
\item[\texttt{defaults write TeXShop syncWithRedOvals YES}] \hfill \\
When syncing from the source to the Preview window, circle the text in red instead of highlighting it.
\item[\texttt{defaults write TeXShop ShowSyncMarks YES}] \hfill \\
Start with the Sync Marks tool on. Marks are shown only for documents typeset with the old pdfsync package.
\end{description}

{\em External Editors and PDF Updates}
\begin{description}
\item[\texttt{defaults write TeXShop AutoReloadExternalPDF NO}] \hfill \\
With Automatic Preview Update on, do not reload a PDF changed by another program.
\item[\texttt{defaults write TeXShop PdfFileRefresh NO}] \hfill \\
Do not update a PDF file opened by itself in TeXShop when it changes. The PDF of a TeX document is still updated.
\item[\texttt{defaults write TeXShop RefreshTime 2}] \hfill \\
How often, in seconds, TeXShop checks whether a PDF has changed; the default is 1 and the minimum 0.5.
\item[\texttt{defaults write TeXShop BringPdfFrontOnAutomaticUpdate NO}] \hfill \\
When an external editor is used, do not bring the Preview window to the front when the PDF is updated.
\item[\texttt{defaults write TeXShop ExternalEditorTypesetAtStart NO}] \hfill \\
When an external editor is used, do not typeset a document when it is opened, even if its PDF is missing or older than the source.
\item[\texttt{defaults write TeXShop TextMateSync YES}] \hfill \\
When an external editor is used, sync from the Preview window and go to errors in TextMate, using /usr/local/bin/mate.
\item[\texttt{defaults write TeXShop OtherEditorSync YES}] \hfill \\
The same with any editor, using a script /usr/local/bin/othereditor that you provide; it is called with the line number and the file. This takes precedence over TextMateSync.
\end{description}

{\em Windows and Saving}
\begin{description}
\item[\texttt{defaults write TeXShop OpenWithSourceInFront YES}] \hfill \\
When a document is opened, bring the source window in front of the Preview window.
\item[\texttt{defaults write TeXShop TabsAlsoForInputFiles YES}] \hfill \\
Make the magic line useTabs also put the files of \verb|\input| commands in tabs.
\item[\texttt{defaults write TeXShop WatchServer NO}] \hfill \\
Do not watch for changes made by other programs to source files on a server. TeXShop watches only while Autosave and Versions is on in TeXShop Settings.
\item[\texttt{defaults write TeXShop IgnoreStartOnReboot YES}] \hfill \\
Do not reopen TeXShop automatically when you log in, even if it was running when you logged out.
\end{description}

{\em Panels}
\begin{description}
\item[\texttt{defaults write TeXShop matrixsize 20}] \hfill \\
The largest number of rows and columns in the Matrix Panel, from 2 to 100; the default is 12.
\item[\texttt{defaults write TeXShop LPanelOutlines NO}] \hfill \\
Always draw the outlines of the buttons in the LaTeX Panel, not only when the pointer is over them.
\end{description}

\end{document}  